\scanpage{89}
\mspage{88}
des parties $\deleted{\unsure{\widehat F|\widehat H}}$ de $\unsure{\mathcal L|\mathscr L}$
\[
\widehat F=\bigl\{\,\lvert X\rvert\mid X\in\bar F\quad\text{i.e. }X\in\unsure{\mathcal M|\mathscr M},\,X\leq F\,\bigr\}
\]

Ainsi on a
\[
\lvert\overset{\deleted{\unsure{\wedge|{-}}}}{F}\rvert\subset\mathcal L,\qquad
\widehat F\subset\unsure{\mathcal P|P}(L)
\]

et on a
\[
\lvert F\rvert=\deleted{\unsure{\coprod|\mathcal L}}\bigcup_{X\in\bar F}\lvert X\rvert=\bigcup_{A\in\widehat F}A\,.
\]

Cette relation signifie que si $x\unsure{\in|+}\mathcal L$ est tel que $x\ll F$, $\exists\,X\in\bar F$ tel que $x\ll X$, \unsure{\S|(*)} 
% slanted
\margin{Mais pour \unsure{le|la} prouver, on a utilisé \unsure{axiome|l'axiome} \unsure{d'invention|d'induction|d'inversion}}
et pour ceci, voir prop. page \unsure{78|76}. Il y a \unsure{même|\={m}} un plus petit tel $X$ (\unsure{=|:} Cor à la prop).

\uline{Proposition} La multi-ombre $\unsure{\Phi|\mathcal P}=\widehat F$ d'une figure est une famille \uline{admissible} de parties de $\mathcal L$, i.e.
\[
\begin{aligned}
a)\quad &\forall\,\deleted{\unsure{a\in|x\in|A\in}}A\in\Phi,\ \text{posant }\ \unsure{\delta|\partial}A=\bigcup_{\substack{B\subset A\\B\neq A}}B\\
&\text{on a}\qquad \delta A\neq A\quad\text{\unsure{i.e|i.e.}}\quad A^{\unsure{\circ|o}}\overset{\mathrm{def}}{=}A\setminus\delta A\neq\varnothing
\end{aligned}
\]

\[
b)\quad \forall\,A,B\in\Phi,\ \text{on a}\qquad A\cap B=\bigcup_{\substack{C\in\Phi\\C\subset A\cap B}}C
\]
